\documentclass[11pt,a4paper]{article}
\usepackage[margin=23mm]{geometry}
\usepackage{amsmath,amssymb,amsthm}
\setlength{\emergencystretch}{2em}
\newtheorem{theorem}{Theorem}
\newtheorem{corollary}[theorem]{Corollary}
\newcommand{\gen}[1]{\langle#1\rangle}
\title{An assigned equality with a fixed anchor}
\author{Human Compression project: L044 continuation}
\date{}
\begin{document}\maketitle
Let $G=\gen{\rho_0,\ldots,\rho_{r-1}}$ be an ambient
string action, put $H=G_{\ge k}$ and $K=G_{>k}$, and
let $V$ be an actual $H$-orbit. Eligibility refers throughout
to the full ambient omissions $G_j$.

\begin{theorem}
Suppose $x\in V$ is fixed by $K$. Let
$F\subseteq\{k,\ldots,r-1\}$ contain $k$, and suppose
every colour in $F$ has two ambient-eligible edges in $V$.
If $|V|=2|F|$, then $F$ is one initial interval
$\{k,\ldots,a-1\}$ for some $a$.
\end{theorem}
\begin{proof}
Ambient eligibility implies eligibility for the full local
omission of $H|_V$, by subgroup inclusion. The complete
selected-colour equality theorem therefore makes $G_F$
transitive on $V$.

Split $F$ into its maximal consecutive intervals.
Their generated groups commute pairwise, since labels in
different intervals differ by at least two. Every component
except the one containing $k$ lies in $K$ and fixes $x$.
In a transitive product of commuting permutation groups,
a factor fixing $x$ fixes every point: write an element
as $ab$, with $a$ in that factor and $b$ in the product
of the others. Then $x^{ab}=x^b$, and the whole factor
fixes $x^b$ by commutation with $b$.
Thus every later component fixes all of $V$. But every
one of its colours has two actual edges in $V$, a
contradiction. Only the initial component remains.
\end{proof}

\begin{corollary}
The entire three-orbit budget of the acquired theorem
\emph{A four-point internal orbit gives a nonterminal
separator} is impossible.
\end{corollary}
\begin{proof}
Retain every literal hypothesis of that theorem, including
the natural alternating tail, full IP, $n=2r+1$, seven
$H=G_{\ge5}$-fixed points, nontrivial orbits $E,W,V$
of sizes $2,4,|V|$, the specified action of $\delta$,
and intransitivity of $K=G_{>5}$ on $W$.
Its assignments are
\[
 F_E=\varnothing,\quad F_W=\{j\},\quad
 F_V=\{5,\ldots,r-1\}\setminus\{j\},\qquad j\ge6.
\]
The same actual head argument supplies $x\in V$ fixed
by $K$. Moreover
\[
 |V|=n-7-2-4=2r-12=2|F_V|.
\]
The theorem just proved makes $F_V$ an initial interval,
so its only missing tail label $j$ must be $r-1$.
The complete acquired four-point-row theorem, with
exactly these hypotheses, gives $j<r-1$.
This contradiction excludes the budget in full.
\end{proof}

\paragraph{Attribution and remaining scope.}
The selected-colour equality and complete four-point-row
proof are acquired inputs, with every predecessor retained.
The new observation is the action of commuting marked
components on a fixed anchor. It uses no projected IP,
new computation or rank induction. In the remaining
nonfirst internal deficit-two budget it forces the first
assigned set to be an initial interval and the other
assigned set to be terminal; it does not exclude that
budget. Split-orbit pairs, global attachment and full
L044 recovery remain. No formal, original, question or
numerical credit is claimed.
\end{document}
